\begin{abstract}
The Census Bureau counts college students at their campus address under the usual-residence rule. In 2020, approximately 8.6 million students were enumerated at dormitories, off-campus university housing, and other collegiate group-quarters facilities rather than at their pre-enrollment home addresses. This campus-counting convention inflates the population of college towns by 10--35\%, creating systematic advantages for the congressional districts that contain major research universities while deflating the suburban and rural districts that are the true residential communities of those students.

We document the campus-inflation effect across the 50 largest college towns in the United States, with extended case studies for Ithaca (New York), Ann Arbor (Michigan), and Champaign-Urbana (Illinois). We then analyze the partisan direction of the distortion: college towns are among the most Democratic-voting localities in the United States, and counting students at campus systematically inflates Democratic-leaning college-district populations at the expense of the suburban and exurban districts from which students typically originate---districts that lean Republican.

We estimate that approximately 180 census tracts nationally would receive a different district assignment if student populations were reallocated to home addresses for the 2020 redistricting cycle. This is a tract-assignment point estimate: the final paper should report the confidence band implied by the home-address sensitivity run, because off-campus housing and tract-boundary proximity can move a small number of tracts across the threshold. The BISECT redistricting pipeline's \texttt{--population-source} flag can implement either the Census-default campus count or a hypothetical home-address reallocation, permitting systematic sensitivity analysis. Litigation in Maine and New York has directly contested the campus-counting rule; we review those proceedings and assess the constitutional and statutory arguments on both sides. We conclude that while the usual-residence rule is legally permissible, college-town inflation represents a structurally partisan, algorithmically correctable feature of census-based redistricting that warrants greater legislative attention.
\end{abstract}
